% Options for packages loaded elsewhere
\PassOptionsToPackage{unicode}{hyperref}
\PassOptionsToPackage{hyphens}{url}
%
\documentclass[
  14pt,
]{extarticle}
\usepackage{lmodern}
\usepackage{amssymb,amsmath}
\usepackage{ifxetex,ifluatex}
\ifnum 0\ifxetex 1\fi\ifluatex 1\fi=0 % if pdftex
  \usepackage[T1]{fontenc}
  \usepackage[utf8]{inputenc}
  \usepackage{textcomp} % provide euro and other symbols
\else % if luatex or xetex
  \usepackage{unicode-math}
  \defaultfontfeatures{Scale=MatchLowercase}
  \defaultfontfeatures[\rmfamily]{Ligatures=TeX,Scale=1}
  \setmainfont[]{Droid Sans Fallback}
\fi
% Use upquote if available, for straight quotes in verbatim environments
\IfFileExists{upquote.sty}{\usepackage{upquote}}{}
\IfFileExists{microtype.sty}{% use microtype if available
  \usepackage[]{microtype}
  \UseMicrotypeSet[protrusion]{basicmath} % disable protrusion for tt fonts
}{}
\makeatletter
\@ifundefined{KOMAClassName}{% if non-KOMA class
  \IfFileExists{parskip.sty}{%
    \usepackage{parskip}
  }{% else
    \setlength{\parindent}{0pt}
    \setlength{\parskip}{6pt plus 2pt minus 1pt}}
}{% if KOMA class
  \KOMAoptions{parskip=half}}
\makeatother
\usepackage{xcolor}
\IfFileExists{xurl.sty}{\usepackage{xurl}}{} % add URL line breaks if available
\IfFileExists{bookmark.sty}{\usepackage{bookmark}}{\usepackage{hyperref}}
\hypersetup{
  hidelinks,
  pdfcreator={LaTeX via pandoc}}
\urlstyle{same} % disable monospaced font for URLs
\usepackage[margin=1in]{geometry}
\usepackage{longtable,booktabs}
% Correct order of tables after \paragraph or \subparagraph
\usepackage{etoolbox}
\makeatletter
\patchcmd\longtable{\par}{\if@noskipsec\mbox{}\fi\par}{}{}
\makeatother
% Allow footnotes in longtable head/foot
\IfFileExists{footnotehyper.sty}{\usepackage{footnotehyper}}{\usepackage{footnote}}
\makesavenoteenv{longtable}
\setlength{\emergencystretch}{3em} % prevent overfull lines
\providecommand{\tightlist}{%
  \setlength{\itemsep}{0pt}\setlength{\parskip}{0pt}}
\setcounter{secnumdepth}{-\maxdimen} % remove section numbering

\author{}
\date{}

\begin{document}

\hypertarget{week-07-ux4f5cux696dux56deux994b}{%
\section{Week 07 作業回饋}\label{week-07-ux4f5cux696dux56deux994b}}

\begin{quote}
\textbf{評分標準：} 依據 \texttt{week07\_rubric.xlsx}
之六項評分指標，基礎分 100 分，Bonus 最高 +10 分，總分上限為 100。
各大項配分：Q1 Data Simulation (25) ／ Q2 Data Cleaning (20) ／ Q3
Descriptive Stats (20) ／ Q4 Multi-Panel Figure (35) ／ Bonus Panel E
(+5) ／ Bonus Paired t-test (+5)
\end{quote}

\begin{center}\rule{0.5\linewidth}{0.5pt}\end{center}

\hypertarget{ux4f55ux5b98ux81fb-95100}{%
\subsection{114825002 何官臻 ---
95／100}\label{ux4f55ux5b98ux81fb-95100}}

\hypertarget{ux7e3dux8a55}{%
\subsubsection{總評}\label{ux7e3dux8a55}}

四道基礎題完成度完整、資料結構與命名皆符合規範，2×2
圖表版面乾淨俐落，Panel B 數值標註、Panel C 組平均虛線、Panel D
speed-accuracy
散佈圖皆達標。程式整體簡潔、每段都有清楚的中文註解。本次未嘗試
Bonus（Panel E 練習/疲勞曲線、paired
t-test），因此最終分數落在基礎分上限；若能補上其中一項 Bonus 則容易達
100。程式碼中有少數註解拼字錯誤（\texttt{tramsform}、\texttt{impulsiable}、第
0 cell 的 \texttt{matplotlb}），不影響執行也不扣分，但建議下次提交前快速
spell-check。

\hypertarget{ux5404ux9805ux8a55ux8a9e}{%
\subsubsection{各項評語}\label{ux5404ux9805ux8a55ux8a9e}}

\begin{longtable}[]{@{}lll@{}}
\toprule
\begin{minipage}[b]{0.10\columnwidth}\raggedright
指標\strut
\end{minipage} & \begin{minipage}[b]{0.04\columnwidth}\raggedright
得分（Base）\strut
\end{minipage} & \begin{minipage}[b]{0.78\columnwidth}\raggedright
評語\strut
\end{minipage}\tabularnewline
\midrule
\endhead
\begin{minipage}[t]{0.10\columnwidth}\raggedright
Q1 Data Simulation (NumPy)\strut
\end{minipage} & \begin{minipage}[t]{0.04\columnwidth}\raggedright
25/25\strut
\end{minipage} & \begin{minipage}[t]{0.78\columnwidth}\raggedright
\texttt{np.random.default\_rng(2026)} ✓，以
\texttt{subjects\ =\ {[}f"P\{i:02d\}"\ for\ i\ in\ range(1,\ 21){]}}
產生 ID 列表後，使用巢狀迴圈依 subject × condition 填入資料；分配參數
Normal(520,80)／Normal(620,100)、p=0.95／0.85 全部正確。欄位命名
\texttt{subject}／\texttt{trial}／\texttt{condition}／\texttt{rt\_ms}／\texttt{correct}
符合規範。shape 與 head(10)、dtypes 皆印出。\strut
\end{minipage}\tabularnewline
\begin{minipage}[t]{0.10\columnwidth}\raggedright
Q2 Data Cleaning\strut
\end{minipage} & \begin{minipage}[t]{0.04\columnwidth}\raggedright
20/20\strut
\end{minipage} & \begin{minipage}[t]{0.78\columnwidth}\raggedright
三步驟（errors → per-subject outlier via \texttt{groupby.transform} →
implausible RT）依序套用，每步均以「原始 −
已扣除」公式計算移除數量（\texttt{outlier\_removed\ =\ (original\_count\ -\ error\_removed)\ -\ len(df\_clean)}），邏輯正確；CSV
輸出 ✓。小註：註解有拼字 \texttt{tramsform}、\texttt{impulsiable}，建議
spell-check。\strut
\end{minipage}\tabularnewline
\begin{minipage}[t]{0.10\columnwidth}\raggedright
Q3 Descriptive Stats (groupby)\strut
\end{minipage} & \begin{minipage}[t]{0.04\columnwidth}\raggedright
20/20\strut
\end{minipage} & \begin{minipage}[t]{0.78\columnwidth}\raggedright
\texttt{groupby({[}\textquotesingle{}subject\textquotesingle{},\textquotesingle{}condition\textquotesingle{}{]}){[}\textquotesingle{}rt\_ms\textquotesingle{}{]}.agg({[}\textquotesingle{}mean\textquotesingle{},\textquotesingle{}std\textquotesingle{},\textquotesingle{}median\textquotesingle{},\textquotesingle{}count\textquotesingle{}{]})}
✓；\texttt{pivot\_table} 重塑後直接 \texttt{incongruent\ -\ congruent}
得 interference cost ✓；利用
\texttt{pivot\_df{[}{[}\textquotesingle{}congruent\textquotesingle{},\textquotesingle{}incongruent\textquotesingle{}{]}{]}.agg({[}\textquotesingle{}mean\textquotesingle{},\textquotesingle{}sem\textquotesingle{}{]})}
一次算出 between-subjects SEM，寫法簡潔；SEM vs SD 的使用正確，Range
輸出也齊全。\strut
\end{minipage}\tabularnewline
\begin{minipage}[t]{0.10\columnwidth}\raggedright
Q4 Multi-Panel Figure\strut
\end{minipage} & \begin{minipage}[t]{0.04\columnwidth}\raggedright
30/35\strut
\end{minipage} & \begin{minipage}[t]{0.78\columnwidth}\raggedright
2×2 版面 ✓，suptitle 題文完全符合
\texttt{"Stroop\ Task:\ Experiment\ Summary"} ✓，Panel A histogram +
平均虛線 ✓、Panel B 長條加 SEM errorbar
且文字註解為平均值（\texttt{f\textquotesingle{}\{height:.1f\}\textquotesingle{}}）✓。可再改進之處：（1）Panel
C 使用 \texttt{sort\_values(ascending=True)} 未加
\texttt{invert\_yaxis()}，實際呈現是最小的 cost 在最上方，與題目「sorted
from largest to smallest」略有方向差（−1）；（2）Panel D 使用原始
\texttt{df} 計算 accuracy，但 RT 使用 \texttt{pivot\_df}
的清理後均值，兩者資料來源不一致，建議都以 \texttt{df\_clean}
為依據；再者 Panel D 沒有呼叫 \texttt{plt.tight\_layout} 之後就
savefig，但實務上圖沒有溢出（−2）；（3）沒有嘗試 Bonus Panel E 與
t-test（不在 Q4 扣分範圍，已反映在後段）。整體而言 Panel B、C
的配色與數值標註呈現清楚、可讀性佳。\strut
\end{minipage}\tabularnewline
\begin{minipage}[t]{0.10\columnwidth}\raggedright
\textbf{Base 小計}\strut
\end{minipage} & \begin{minipage}[t]{0.04\columnwidth}\raggedright
\textbf{95/100}\strut
\end{minipage} & \begin{minipage}[t]{0.78\columnwidth}\raggedright
\strut
\end{minipage}\tabularnewline
\begin{minipage}[t]{0.10\columnwidth}\raggedright
Bonus Panel E (+5)\strut
\end{minipage} & \begin{minipage}[t]{0.04\columnwidth}\raggedright
\textbf{+0}\strut
\end{minipage} & \begin{minipage}[t]{0.78\columnwidth}\raggedright
未嘗試。Panel E 是分析「練習效應／疲勞效應」的好機會，只需
\texttt{df\_clean.groupby({[}\textquotesingle{}trial\textquotesingle{},\textquotesingle{}condition\textquotesingle{}{]}){[}\textquotesingle{}rt\_ms\textquotesingle{}{]}.mean().unstack()}
再用 \texttt{ax.plot()} 兩條線即可。\strut
\end{minipage}\tabularnewline
\begin{minipage}[t]{0.10\columnwidth}\raggedright
Bonus Paired t-test (+5)\strut
\end{minipage} & \begin{minipage}[t]{0.04\columnwidth}\raggedright
\textbf{+0}\strut
\end{minipage} & \begin{minipage}[t]{0.78\columnwidth}\raggedright
未嘗試。\texttt{scipy.stats.ttest\_rel(pivot\_df{[}\textquotesingle{}incongruent\textquotesingle{}{]},\ pivot\_df{[}\textquotesingle{}congruent\textquotesingle{}{]})}
一行即可印出 t／p 值，是後續 Week 08／期中作業統計分析的重要基礎。\strut
\end{minipage}\tabularnewline
\begin{minipage}[t]{0.10\columnwidth}\raggedright
\textbf{最終分數}\strut
\end{minipage} & \begin{minipage}[t]{0.04\columnwidth}\raggedright
\textbf{95/100}\strut
\end{minipage} & \begin{minipage}[t]{0.78\columnwidth}\raggedright
\strut
\end{minipage}\tabularnewline
\bottomrule
\end{longtable}

\hypertarget{ux5f8cux7e8cux5efaux8b70}{%
\subsubsection{後續建議}\label{ux5f8cux7e8cux5efaux8b70}}

\begin{enumerate}
\def\labelenumi{\arabic{enumi}.}
\tightlist
\item
  \textbf{補 Bonus 的思維習慣}：Bonus
  題目多半是把已經算出的資料換個角度再呈現（或跑一個統計），不需要重建
  pipeline。下週開始可以練習：每完成一題後，想想「能不能再從這份資料多問一個問題」，這是
  big-data 分析的基本姿態。
\item
  \textbf{排序方向的細節}：在橫條圖 (\texttt{barh}) 用
  \texttt{invert\_yaxis()}
  讓最大值置頂是科學圖常見慣例，讀者視線自然由上而下閱讀；單純
  \texttt{sort\_values(ascending=True)} 會讓最小值置頂。
\item
  \textbf{資料來源一致性}：Panel D 的 RT 與 accuracy 若一個來自
  \texttt{df\_clean} 另一個來自 \texttt{df}，雖然差異很小（clean 後
  accuracy 永遠是 1），但在 log 與 debug 時容易產生混亂；建議以
  \texttt{df\_clean} + 原始 \texttt{df} 的 accuracy
  做明確的欄位合併（\texttt{merge}）。
\item
  \textbf{拼字／第 0 cell}：\texttt{\%pip\ install\ matplotlb}
  實際上會跑出 PyPI 找不到套件的錯誤（應為 \texttt{matplotlib}），顯示
  Colab 環境已預裝才沒影響 notebook；下次建議使用
  \texttt{\%pip\ install\ -q\ matplotlib} 以降低噪音。
\end{enumerate}

\end{document}
